% W5i84_HardProblems.tex
% DFD 20-Agent Research Programme --- Iteration 84
% Hard Problems, Honest Negatives, and Open Questions

\documentclass[11pt,a4paper]{article}
\usepackage[utf8]{inputenc}
\usepackage{amsmath,amssymb}
\usepackage{booktabs}
\usepackage{geometry}
\usepackage{xcolor}
\usepackage{enumitem}
\geometry{margin=2.5cm}

\title{W5i84: Hard Problems and Honest Negatives\\[6pt]
\large Baryonic Feedback Limitation $\cdot$ Multi-Probe vs.\ Single-Probe Detection $\cdot$ Submission Timing}
\author{DFD 20-Agent Research Programme}
\date{2026-03-24}

\begin{document}
\maketitle

\section{Honest Negatives: One Update}

\textbf{HN-14 (New)}: The combined Euclid multi-probe DFD detection is $4.5\sigma$, below the conventional $5\sigma$ discovery threshold. Adding CMB-S4 pushes to $5.5\sigma$, but this requires both experiments to deliver design sensitivity. \textbf{Severity: INFORMATIONAL.}

Total honest negatives: \textbf{14}.

\section{Hard Problem: Baryonic Feedback in N-Body}

This is the most significant systematic limitation of the N-body programme.

\textbf{The problem}: At $k > 1\,h\,\text{Mpc}^{-1}$, baryonic effects (AGN feedback, gas cooling, star formation) modify the matter power spectrum by $2$--$5\%$. The DFD signal at these scales ($\sim 2.5\%$) is comparable. DFD N-body simulations are dark-matter-only and cannot separate these effects.

\textbf{Quantitative assessment}:
\begin{center}
\begin{tabular}{lccc}
\toprule
$k$ [$h$ Mpc$^{-1}$] & DFD Signal & Baryonic Effect & Signal/Systematic \\
\midrule
0.1 & $1.1\%$ & $< 0.1\%$ & $> 10$ \\
0.3 & $1.5\%$ & $0.5\%$ & 3 \\
0.5 & $1.9\%$ & $1.5\%$ & 1.3 \\
1.0 & $2.5\%$ & $3.0\%$ & 0.8 \\
2.0 & $3.0\%$ & $5.0\%$ & 0.6 \\
\bottomrule
\end{tabular}
\end{center}

\textbf{DFD results are only clean at $k < 0.5\,h\,\text{Mpc}^{-1}$} where the signal exceeds the baryonic uncertainty. At $k > 1$, baryonic effects dominate and the DFD signal cannot be isolated without hydrodynamic simulations.

\textbf{Mitigation}:
\begin{enumerate}[nosep]
\item N-body paper I will present results at $k < 0.5$ as primary findings
\item Baryonification models (BCM, Schneider+2019) will be applied at $0.5 < k < 3$
\item The $k > 1$ regime will be flagged as ``DFD + baryon model'' predictions
\item Long-term: DFD-AREPO or DFD-GADGET4 hydro simulations (post-programme)
\end{enumerate}

\section{Hard Problem: Multi-Probe Requirement}

Finding 206 shows that no single Euclid probe exceeds $3\sigma$ for DFD detection. This has strategic implications:

\begin{itemize}[nosep]
\item \textbf{Positive}: DFD's predictions form a consistent \emph{pattern} across many probes. The $4.5\sigma$ combined detection relies on the same underlying parameter ($\psi_0$) being detected in independent channels. This is scientifically compelling.

\item \textbf{Negative}: Referees may argue that ``no single $5\sigma$ detection'' means DFD is untestable. This is incorrect (the combined significance is near-$5\sigma$) but will require careful communication.

\item \textbf{Analogy}: The Higgs discovery combined multiple channels (diphoton, 4-lepton, WW, etc.), none exceeding $5\sigma$ individually but combining to $> 5\sigma$. DFD's situation is analogous.
\end{itemize}

\textbf{Communication strategy}: The Euclid pre-registration paper must frame the multi-probe combination as the primary test, with individual probes as supporting evidence. The analogy to multi-channel particle physics discoveries should be explicit.

\section{Hard Problem: Submission Timing}

The programme has 11 papers at 100\% and 1 submitted. The decision to begin Phase 1 submission (PRL + Physics Reports) has been repeatedly deferred. Analysis of the timing:

\textbf{Arguments for immediate submission}:
\begin{enumerate}[nosep]
\item All scientific prerequisites met ($c_T$, $w_a$ resolved)
\item Euclid pre-registration requires prior publication of core theory
\item First-mover advantage in the ``testable modified gravity'' space
\item The $\alpha$ PRL result is time-independent; no benefit to waiting
\end{enumerate}

\textbf{Arguments for delay}:
\begin{enumerate}[nosep]
\item The BD letter (now complete) strengthens the PRL submission
\item The $c_T$ result (near-complete) pre-empts a referee objection
\item Simultaneous submission of PRL + BD + $c_T$ creates a ``package''
\end{enumerate}

\textbf{Recommendation}: Submit at i85. The PRL + BD letter + $c_T$ paper form a coherent first-contact package. Physics Reports can follow immediately.

\section{Hard Problem: $\Sigma m_\nu$ Distinguishability}

The i84 verification revealed that CMB-S4's ability to distinguish $\Sigma m_\nu = 61.4$ meV (DFD) from $\Sigma m_\nu = 58$ meV (minimum normal hierarchy) is only $0.2\sigma$. This means:

\begin{itemize}[nosep]
\item CMB-S4 can detect $\Sigma m_\nu \neq 0$ at $4.1\sigma$
\item But it \emph{cannot} distinguish DFD's specific prediction from the minimum
\item JUNO ($\sigma = 10$ meV for individual $m_i$) is needed for the $3.4$ meV distinction
\end{itemize}

This weakens Y1 as a near-term test. DFD's $\Sigma m_\nu$ prediction is testable, but the discriminating power between DFD and ``generic normal hierarchy'' is limited to JUNO.

\section{Risk Register: i84 Update}

\begin{center}
\begin{tabular}{lccc}
\toprule
\textbf{Risk} & \textbf{i82} & \textbf{i84} & \textbf{Trend} \\
\midrule
$w_a \neq 0$ confirmed & LOW & LOW & Stable \\
$E_G$ undetectable & LOW & LOW & Stable \\
PRL rejection & MODERATE & MODERATE & Stable \\
$S_8$ mismatch & LOW & LOW & Stable \\
N-body baryonic contamination & LOW & \textbf{MODERATE} & $\uparrow$ Identified \\
Multi-probe below $5\sigma$ & --- & \textbf{LOW} & New \\
Submission delay & --- & \textbf{LOW} & New \\
\bottomrule
\end{tabular}
\end{center}

The baryonic contamination risk is elevated to MODERATE because it limits the N-body paper's $k$-range. All other risks remain stable or low.

\section{Open Questions for i85 (3/4 Mark Assessment)}

\begin{enumerate}[nosep]
\item What has the programme achieved in i52--i85?
\item What is the realistic target for i86--i100?
\item Are the 3 YELLOW items genuinely experiment-limited?
\item What is the optimal submission sequence?
\item Can the programme reach $5\sigma$ combined detection without CMB-S4?
\item What are the top 3 risks for the final quarter?
\end{enumerate}

\end{document}
